\section{Method}
\label{sec:method}

\todofig{overview}{Overview. Frames and their monocular depth maps place planar surfels. Each entity keeps
two information records on the same three placement coordinates, one from the depth prior and one from the
pixels through the differentiable renderer. Their disagreement is tested, and the calibrated records are
fused into a map with an information field.}

\subsection{Setting}
\label{sec:method:setting}
The input is a posed aerial video in strips, with frames $i=1,\dots,N$, images $\mathbf c_i$ and one
monocular depth map $\hat d_i$ per frame from a released model. Depth is related to the true depth $d$ by
$\hat d_i(u)=a_i\,d+b_i+\varepsilon_i(u)$ with unknown scale $a_i$ and offset $b_i$ per frame. The scene is a
set of planar Gaussian surfels \cite{gs_2dgs} with centre $\boldsymbol\mu_k$, tangents
$\mathbf t_{1},\mathbf t_{2}$, normal $\mathbf n$, scales and opacity, rendered with compositing weights
$w_{uk}$ for pixel $u$ and surfel $k$ \cite{gs_gsplat}; rendered depth and colour are
$D_i(u)=\sum_k w_{uk}d_{uk}$ and $\hat{\mathbf c}_i(u)=\sum_k w_{uk}\mathbf a_k$.

An \emph{entity} $m$ is a set of surfels that move together: here the surfels of one ground tile of
\ph{8}~m side with a common normal class. Its \emph{placement} $\boldsymbol\theta_m=(\nu,\tau_1,\tau_2)$ is
an offset along the normal and two tilts about the entity centre $\bar{\boldsymbol\mu}_m$. For a ray with
direction $\mathbf w_u$ that meets surfel $k$ at $\mathbf x$, the depth changes with the placement by the row
\begin{equation}
\mathbf j_{uk}=\frac{\partial d_{uk}}{\partial\boldsymbol\theta_m}
=\frac{1}{\mathbf n^{\top}\mathbf w_u}\bigl[\,1,\ -\mathbf t_1^{\top}(\mathbf x-\bar{\boldsymbol\mu}_m),\
-\mathbf t_2^{\top}(\mathbf x-\bar{\boldsymbol\mu}_m)\bigr].
\label{eq:depthjac}
\end{equation}

\subsection{Two Records on One Placement}
\label{sec:method:records}
\runin{Depth record.} Frame $i$ contributes to entity $m$ the information and pull
\begin{align}
\Lambda^{D}_{mi}&=\sum_{u}\sum_{k\in m}\frac{w_{uk}\,a_i^{2}}{\sigma_i^{2}}\,\mathbf j_{uk}^{\top}\mathbf j_{uk},
\label{eq:depthrecord}\\
\boldsymbol\eta^{D}_{mi}&=\sum_{u}\sum_{k\in m}\frac{w_{uk}\,a_i}{\sigma_i^{2}}\,\mathbf j_{uk}^{\top}r_i(u),
\nonumber
\end{align}
with $r_i$ the aligned depth residual and $\sigma_i$ the depth noise scale of the frame. The weights enter
linearly (lumped information), which majorises the true Hessian, and $(a_i,b_i)$ are marginalised by a
Schur complement, so a frame says nothing in the directions its own scale and offset can explain.

\runin{Image record.} Moving surfel $k$ changes the image in two steps. A change of its kernel changes its
own weight and, through the transmittance, the weights of everything behind it, so the image responds with
the \emph{contrast} $w_{uk}(\mathbf a_k-\mathbf b_{uk})$, where $\mathbf b_{uk}$ is the colour composited
behind $k$. Moving the plane along the ray \emph{slides} the intersection in the tangent frame by
$\mathbf m_{uk}$, the tangential component of the ray in units of the surfel scales. With
$\mathbf s_{uk}$ the intersection in the same units,
\begin{equation}
\frac{\partial\hat{\mathbf c}_i(u)}{\partial\boldsymbol\theta_m}=\mathbf g_{uk}\,\mathbf j_{uk},\qquad
\mathbf g_{uk}=-\,w_{uk}\,(\mathbf a_k-\mathbf b_{uk})\,(\mathbf s_{uk}^{\top}\mathbf m_{uk}).
\label{eq:factor}
\end{equation}
The image Jacobian factors through the depth Jacobian: both sources act on the same row $\mathbf j_{uk}$
and differ by a scalar weight per pixel and surfel. The identity is first-order, with surfel order and the
colour behind held fixed, and holds because a surfel is a plane; it is checked by finite differences on a
reference renderer (Section~\ref{sec:setup}). The image record is
\begin{equation}
\Lambda^{I}_{mi}=\sum_{u}\sum_{k\in m}\frac{\lVert\mathbf g_{uk}\rVert^{2}}{\sigma_I^{2}}\,
\mathbf j_{uk}^{\top}\mathbf j_{uk},
\label{eq:imagerecord}
\end{equation}
with the pull $\boldsymbol\eta^{I}_{mi}$ formed likewise from the colour residual; both come from one
backward pass of the image loss per frame, with visibility and occlusion supplied by the rasteriser. No
features are matched: frames are associated only through the compositing weights of shared surfels.
Records add over frames, $\Lambda^{s}_{m}=\sum_i\Lambda^{s}_{mi}$ for $s\in\{D,I\}$, and over tiles.

\subsection{Honest Information}
\label{sec:method:honest}
\runin{The prior saturates.} Write the placement error of frame $i$'s depth on entity $m$ as
$\boldsymbol\mu_m+\boldsymbol\alpha_{mi}$, a part shared by all frames with covariance $\Sigma_\mu$ and a
part that differs between frames. Fusing $n$ frames leaves the covariance
$\Sigma_\mu+(\Lambda^{D}_{m})^{-1}$, so the honest prior information is
\begin{equation}
\tilde\Lambda^{D}_{m}=\bigl(\Sigma_\mu+(\Lambda^{D}_{m})^{-1}\bigr)^{-1}\ \preceq\ \Sigma_\mu^{-1},
\label{eq:saturation}
\end{equation}
which stops growing with $n$. With $\varrho$ the share of the single-frame error variance that is common,
the fused variance is $\varrho+(1-\varrho)/n$ of the single-frame variance, against the $1/n$ that
\eqref{eq:depthrecord} states. $\Sigma_\mu$ is estimated from the correlation of depth residuals between
frames more than a block of $B$ frames apart along the strip.

\runin{The pixel information is deflated.} Rendering residuals are model misfit and are spatially
correlated, so \eqref{eq:imagerecord} overstates the information. A dispersion $\hat\varphi$ is estimated
from frame replicates: the frames that see an entity are split into halves, and the scatter between the
placements they imply is compared with the scatter their records state. A second, absolute check injects
a known displacement by moving the entity's surfels and re-rendering, and compares the displacement the
record recovers with the one injected; the ratio of stated to realised information is the honesty factor.
Neither step uses a label.

\subsection{Tension and the Score Test}
\label{sec:method:test}
With $\hat{\boldsymbol\theta}^{s}_{m}=(\Lambda^{s}_{m})^{-1}\boldsymbol\eta^{s}_{m}$ the placement each
source implies, the stretch $\mathbf t_m=\hat{\boldsymbol\theta}^{I}_{m}-\hat{\boldsymbol\theta}^{D}_{m}$ is in
metres and radians, and
\begin{equation}
T_m=\mathbf t_m^{\top}\bigl((\tilde\Lambda^{D}_{m})^{-1}+\hat\varphi\,(\Lambda^{I}_{m})^{-1}\bigr)^{-1}\mathbf t_m
\label{eq:tension}
\end{equation}
is $\chi^{2}$ with three degrees of freedom if both sources see one surface; $E_m=T_m/2$ is an energy in
nats. A joint fit would hide a fault, because appearance can be repainted to agree with a wrong surface.
The test is therefore a score test: the surface stays where the depth prior placed it, the appearance is
taken projectively from a set of source frames, and $\boldsymbol\eta^{I}_{m}$, $\Lambda^{I}_{m}$ are
accumulated on disjoint judging frames, so $\mathbf t_m=(\Lambda^{I}_{m})^{-1}\boldsymbol\eta^{I}_{m}$ is
the step the judging pixels ask for. A depth error that every frame repeats does not appear in the
disagreement among depth maps; it appears in $T_m$.

\subsection{Fusion and Scoring}
\label{sec:method:fusion}
The fused record is $\Lambda_m=\tilde\Lambda^{D}_{m}+\hat\varphi^{-1}\Lambda^{I}_{m}$, with the depth
record left out where $T_m$ exceeds the threshold of a stated false-alarm rate. Where
$\hat\varphi^{-1}\Lambda^{I}_{m}$ is negligible the entity is reported as unsupported by pixels. A map is
scored by the \emph{surface log-score}, the negative log density of the true placement of each entity
under the map, taken on the signed error along the normal and averaged over entities. It is a proper
scoring rule \cite{stat_gneiting_scoring}: it cannot be improved by inflating or deflating the stated
information, and it falls when either the geometry or its error bar improves.
